% ************************** Thesis Abstract *****************************
% Use `abstract' as an option in the document class to print only the titlepage and the abstract.
\begin{abstract}

The formation of carbonaceous nanoparticles such as soot and Carbon Black (CB) is a complex phenomenon that involves heat transfer, fluid dynamics, chemical kinetics, and particle dynamics. Understanding the effect of process parameters on the morphology, internal structure, and composition of carbonaceous nanoparticles is essential for assessing the environmental impact of soot as a major air pollutant and for producing specific grades of CB, a key reinforcing ingredient in rubber and tire industries. However, the wide span of time and length scales involved in the process, together with strong temperature and pressure sensitivities, complicates accurate prediction of CB and soot properties.

This project is aimed at developing a robust computational package, titled \textit{Omnisoot}, to describe the formation of carbonaceous nanoparticles, such as soot and CB, from the reactions of gaseous hydrocarbons. Omnisoot extends the chemical-kinetics capabilities of Cantera for gas-phase chemistry to the simulation of solid particles based on well-justified assumptions. It integrates constant volume, imposed pressure, perfectly stirred, and plug flow reactor models with three widely used inception models, as well as a new inception model based on E-Bridge bond formation. The particle dynamics are described using two population balance models, a monodisperse model and a fixed sectional model, that account for the evolving fractal-like structure of soot agglomerates. The implementation captures soot inception, surface growth, and oxidation coupled with detailed gas-phase chemistry to close the mass and energy balances of the gas–particle system. A step-by-step validation approach was followed to ensure the reliability of all sub-models, enabling the package to serve as an integrated process design tool to predict soot mass, morphology, and composition under varying process conditions.

The developed tool was then used to investigate the temperature dependence of the implemented inception models and its impact on predicted soot mass, morphology, and size distribution for various use-cases. The importance of combined soot yield and morphology data as calibration targets was emphasized for constraining the inception flux and surface growth rate. Special attention was paid to methane pyrolysis because of its practical significance in plasma reactors for co-generation of hydrogen and CB.

Shock-tube pyrolysis of methane was studied at high fuel loadings in post-reflected-shock temperature and pressure ranges relevant to CB synthesis. Detailed comparison of model predictions with time-resolved species mole fractions and soot volume fraction ($f_v$) recorded in collaboration with Stanford University enabled evaluation of several kinetic mechanisms and highlighted the major carbon flux pathways from small hydrocarbons to soot precursors. The impact of fuel loading on gas chemistry and soot formation was interpreted in terms of the endothermicity of pyrolysis. Omnisoot successfully predicted the time history of $f_v$ and the final primary particle diameter ($d_p$) in agreement with measurements. The different stages of soot formation were identified by analyzing the time histories of predicted $f_v$, inception flux, and surface growth rates. The evolution of the hydrogen-to-carbon ratio (H/C) predicted using Omnisoot showed strong correspondence with the maturity index ($\beta$) inferred from extinction ratios at multiple wavelengths, highlighting its potential as a quantitative, time-resolved probe of soot maturity.
 
\end{abstract}
